\section{Security Analysis}
\label{sec:security}

We prove that the Accountable Redaction protocol satisfies the four
security properties defined in Section~\ref{sec:model}. Each proof is
structured as a reduction to a known cryptographic primitive.

\subsection{Property 1: Redaction Soundness}
\label{subsec:soundness}

\begin{theorem}
  A computationally bounded adversary $\mathcal{A}$ cannot produce a
  proof $\pi_{\mathit{stat}}$ accepted by $\mathsf{Verify}$ for a
  redaction batch $R$ that violates $\varepsilon$-statistical consistency
  (Definition~\ref{def:eps-stat}), except with negligible probability.
\end{theorem}

\begin{proof}
Proceed by contradiction. Suppose $\mathcal{A}$ outputs a tuple
$(R, \pi_{\mathit{stat}})$ such that $\mathsf{Verify}(\mathit{vk},
\mathbf{x}, \pi_{\mathit{stat}}) = 1$ where $\mathbf{x}$ encodes a
transition with true statistical deviation $\Delta_g > \varepsilon_g$ for
some $g \in \mathcal{G}$.

By the knowledge soundness of the underlying SNARK (Groth16 or Plonk),
there exists a polynomial-time extractor $\mathcal{E}$ such that, given
$\pi_{\mathit{stat}}$, $\mathcal{E}$ produces a witness $w$ satisfying
every constraint in $\mathcal{C}$. In particular, $w$ must satisfy the
consistency constraint:
\[
  \bigl| q_g^{(t)} - q_g^{(t+1)} \bigr| \;\leq\; \varepsilon_g \cdot
  \mathsf{S}
\]
where $q_g^{(t)}, q_g^{(t+1)}$ are the fixed-point quotients derived
from the per-group statistics (Section~\ref{subsec:fixedpoint}). A
witness satisfying this constraint implies the true deviation $\Delta_g
\leq \varepsilon_g$, contradicting the assumption that $\Delta_g >
\varepsilon_g$.

Since the constraint is unsatisfiable under the assumed conditions, no
valid witness exists, and no computationally bounded prover can produce an
accepted proof without breaking the knowledge soundness of the SNARK.
\end{proof}

\subsection{Property 2: Content Indistinguishability}
\label{subsec:privacy}

\begin{theorem}
  For any redacted index $i \in R$, the post-redaction commitment $C_i'$
  is computationally indistinguishable from a fresh null commitment
  $C_{\bot}$. In particular, the demographic attribute $g_i$ and the
  decision content of $\ell_i$ are irrecoverable from the public log
  state.
\end{theorem}

\begin{proof}
The public view after redaction consists of the leaf commitment $C_i' =
C_{\bot}$ (a null Pedersen commitment) and the proof $\pi_{\mathit{stat}}$.
By the zero-knowledge property of the SNARK, $\pi_{\mathit{stat}}$ is
computationally indistinguishable from a simulated proof that does not
use the witness; hence it reveals no information about $g_i$ or $\ell_i$
beyond what is implied by the public inputs.

Suppose $\mathcal{A}$ can distinguish the group $g_i \in \{g_0, g_1\}$
of a redacted entry with non-negligible advantage $\delta$. Construct a
simulator $\mathcal{S}$ against the hiding game for Pedersen commitments:
\begin{enumerate}
  \item The challenger samples $b \gets \{0,1\}$ and provides
  $C^* = \mathsf{Com}(g_b;\, r)$ to $\mathcal{S}$.
  \item $\mathcal{S}$ embeds $C^*$ as the leaf for index $i$ in a
  simulated log and uses the SNARK simulator (whose existence is
  guaranteed by the zero-knowledge property) to produce a valid-looking
  $\pi_{\mathit{stat}}$ without the real witness.
  \item $\mathcal{S}$ presents $(C^*, \pi_{\mathit{stat}})$ to
  $\mathcal{A}$, which outputs $\hat{b}$. $\mathcal{S}$ outputs
  $\hat{b}$.
\end{enumerate}
$\mathcal{S}$'s advantage equals $\mathcal{A}$'s advantage $\delta$ in
breaking the computational hiding of the Pedersen commitment over the
underlying group. Since Pedersen commitments are computationally hiding
under the Discrete Logarithm assumption in the generic group model,
$\delta$ is negligible.
\end{proof}

\subsection{Property 3: Merkle Continuity}
\label{subsec:continuity}

\begin{theorem}
  A valid redaction batch $R$ targeting indices $i \in R$ cannot modify,
  remove, or invalidate the Merkle membership of any entry $j \notin R$.
  Existing consistency proofs computed against $\mathit{rt}_t$ remain
  verifiable against $\mathit{rt}_{t+1}$.
\end{theorem}

\begin{proof}
A redaction replaces leaf $i$ with $C_{\bot}$ and recomputes the path
from leaf $i$ to the root. All sibling nodes on paths not intersecting
with the redacted path are unchanged. Modifying leaf $j \notin R$ while
producing a valid root $\mathit{rt}_{t+1}$ would require finding two
distinct preimages yielding the same internal node, i.e., a collision in
the hash function $H$ (Poseidon). By collision resistance of $H$, this
occurs with at most negligible probability.

Merkle consistency proofs (RFC~9162~\cite{rfc6962}) allow any party
holding $\mathit{rt}_t$ to verify that $\mathit{rt}_{t+1}$ was obtained
by a valid append or leaf substitution without re-downloading the log.
This property is inherited without modification, since nullifying a leaf
is structurally equivalent to updating its value.
\end{proof}

\subsection{Property 4: Operator Non-Frameability}
\label{subsec:nonframe}

\begin{theorem}
  (1) A malicious operator cannot produce a valid $\pi_{\mathit{auth}}$
  for a redaction not requested by the data subject. (2) A malicious
  data subject cannot credibly accuse an honest operator of performing an
  unauthorized redaction that did not occur.
\end{theorem}

\begin{proof}[Proof of Direction~1 (Operator malicious)]
Every accepted redaction requires $\pi_{\mathit{auth}} = (i,
\sigma_{\mathit{del}}, \mathit{pk}_S)$ where
$\mathsf{Verify}_{\mathit{pk}_S}(\sigma_{\mathit{del}},
\texttt{DELETE}\,\|\,i\,\|\,\mathit{ts}) = 1$ and $\mathit{pk}_S$ matches
the identity embedded in $\ell_i$ at insertion time.

Suppose operator $\mathcal{A}$ produces a valid redaction for entry $i$
without the subject's consent. $\mathcal{A}$ must produce a valid
signature $\sigma^*$ on a message $m^* = (\texttt{DELETE}\,\|\,i\,\|\,
\mathit{ts})$ under $\mathit{pk}_S$. Since the subject did not consent,
$m^*$ was never signed by the subject. Construct $\mathcal{S}$ against
EUF-CMA of the signature scheme: $\mathcal{S}$ simulates the protocol
and forwards all legitimate signing queries; when $\mathcal{A}$ outputs
$\sigma^*$, $\mathcal{S}$ outputs $(m^*, \sigma^*)$ as the forgery.
This contradicts EUF-CMA security of the signature scheme.
\end{proof}

\begin{proof}[Proof of Direction~2 (Subject malicious)]
Suppose subject $\mathcal{A}$ claims the operator redacted entry $i$
without consent. Two cases:
\begin{enumerate}
  \item \textbf{The redaction occurred.} Then the Redaction Ledger
  contains an entry for batch $R \ni i$ with $\pi_{\mathit{auth}}^{(i)}$.
  If $\pi_{\mathit{auth}}^{(i)}$ contains a valid signature under
  $\mathit{pk}_S$, the subject did authorise it; the accusation is
  false. Claiming otherwise requires the subject to deny their own valid
  signature, which is possible only if the signature scheme is
  repudiable --- contradicting the non-repudiation guarantee of
  EUF-CMA-secure schemes.
  \item \textbf{The redaction did not occur.} Then the Redaction Ledger
  contains no entry for index $i$. By Property~3 (Merkle Continuity),
  $\ell_i$ is still present in the main log with a valid inclusion proof
  against the current root $\mathit{rt}$. The operator produces this
  proof, refuting the claim. Fabricating a false absence would require
  forging a Merkle exclusion proof, which requires finding a hash
  collision (contradicts collision resistance of $H$).
\end{enumerate}
In both cases, the subject cannot mount a successful false accusation.
\end{proof}
